\documentclass[11pt,a4paper]{article}

\newcommand{\magiversion}{v0.8.2-beta}
% Shared preamble for the MAGI user manual.
\usepackage[T1]{fontenc}
\usepackage[utf8]{inputenc}
\usepackage{lmodern}
\usepackage[a4paper,margin=2.4cm,headheight=15pt]{geometry}
\usepackage{amsmath,amssymb}
\usepackage{graphicx}
\usepackage{booktabs}
\usepackage{longtable}
\usepackage{array}
% [table] pulls in colortbl, needed for the \rowcolor highlighting in the
% per-layer usage table (Visualization section).
\usepackage[table]{xcolor}
\usepackage{listings}
\usepackage{enumitem}
\usepackage{microtype}
\usepackage{caption}
\usepackage{fancyhdr}
\usepackage{titlesec}
\usepackage{tcolorbox}
\usepackage[hidelinks,breaklinks]{hyperref}

\tcbuselibrary{skins,breakable}

% ---------------------------------------------------------------- colours
\definecolor{magiblue}{HTML}{1F4E79}
\definecolor{magigrey}{HTML}{4A4A4A}
\definecolor{codebg}{HTML}{F6F8FA}
\definecolor{codekw}{HTML}{A626A4}
\definecolor{codestr}{HTML}{50A14F}
\definecolor{codecom}{HTML}{9CA0A4}
\definecolor{warnbg}{HTML}{FFF4E5}
\definecolor{warnrule}{HTML}{C77700}
\definecolor{notebg}{HTML}{EEF4FB}

\hypersetup{
  colorlinks=true,
  linkcolor=magiblue,
  urlcolor=magiblue,
  citecolor=magiblue,
}

% ---------------------------------------------------------------- headers
\pagestyle{fancy}
\fancyhf{}
\fancyhead[L]{\small\textcolor{magigrey}{MAGI \magiversion{} --- User Manual}}
\fancyhead[R]{\small\textcolor{magigrey}{\leftmark}}
\fancyfoot[C]{\small\thepage}
\renewcommand{\headrulewidth}{0.4pt}

\titleformat{\section}{\Large\bfseries\color{magiblue}}{\thesection}{0.7em}{}
\titleformat{\subsection}{\large\bfseries\color{magiblue}}{\thesubsection}{0.7em}{}

% ---------------------------------------------------------------- listings
\lstdefinestyle{magi}{
  backgroundcolor=\color{codebg},
  basicstyle=\ttfamily\footnotesize,
  keywordstyle=\color{codekw}\bfseries,
  stringstyle=\color{codestr},
  commentstyle=\color{codecom}\itshape,
  breaklines=true,
  breakatwhitespace=false,
  showstringspaces=false,
  frame=single,
  rulecolor=\color{codebg!60!black!20},
  framesep=5pt,
  xleftmargin=6pt,
  columns=fullflexible,
  keepspaces=true,
  upquote=true,
  literate={-}{{-}}1 {_}{{\_}}1,
}
\lstset{style=magi}
\lstdefinelanguage{shell}{
  keywords={python,pytest,pip,cd,export,bash},
  sensitive=true,
  comment=[l]{\#},
}
\newcommand{\py}[1]{\lstinline[language=Python]|#1|}
\newcommand{\code}[1]{\texttt{\small #1}}
% Long snake_case identifiers inside \code give TeX no break points, which is
% what runs a line into the margin. \cbrk marks a legal break inside one
% (no hyphen inserted); a little \emergencystretch absorbs the rest.
\newcommand{\cbrk}{\discretionary{}{}{}}
\setlength{\emergencystretch}{1.2em}

% ---------------------------------------------------------------- callouts
\newtcolorbox{magiwarn}[1][]{
  breakable, enhanced, colback=warnbg, colframe=warnrule,
  boxrule=0pt, leftrule=3pt, arc=1pt,
  fonttitle=\bfseries, coltitle=warnrule, title={#1},
  left=7pt, right=7pt, top=5pt, bottom=5pt,
}
\newtcolorbox{maginote}[1][]{
  breakable, enhanced, colback=notebg, colframe=magiblue,
  boxrule=0pt, leftrule=3pt, arc=1pt,
  fonttitle=\bfseries, coltitle=magiblue, title={#1},
  left=7pt, right=7pt, top=5pt, bottom=5pt,
}


\title{\vspace{-1.5cm}\textcolor{magiblue}{\Huge\bfseries MAGI}\\[0.3cm]
  {\Large A CVAE-based generative particle source for Geant4}\\[0.5cm]
  {\large User Manual --- \magiversion}}
\author{Francesco Monastra\\\small INAF --- \texttt{francesco.monastra@inaf.it}}
\date{\today}

\begin{document}
\maketitle
\thispagestyle{empty}

\begin{magiwarn}[This release is a beta]
MAGI \magiversion{} reproduces the \emph{joint} phase-space distribution
(energy--position--direction correlations) and the energy continuum to within
its stated acceptance bars, on three seeds and two reference sources. It does
\textbf{not} yet reproduce spectral-line \emph{intensities}: measured per-line
recovery runs from $0.7\times$ to $4.9\times$ the true value and is unstable
across seeds. Line \emph{positions} are reliable to $\leq 0.5$~eV.

Read Section~\ref{sec:envelope} before quoting any number derived from a
MAGI-generated source. Do not use this release to estimate a fluorescence or
annihilation line flux.
\end{magiwarn}

\vspace{0.4cm}
\tableofcontents
\newpage

% =====================================================================
\section{Introduction}
\label{sec:intro}

\subsection{What MAGI is}

MAGI is a Conditional Variational Autoencoder
(CVAE)~\cite{Kingma2014VAE,Sohn2015CVAE}, implemented in Keras/TensorFlow,
that learns the phase-space distribution of particles crossing a detector
surface in a Geant4~\cite{Agostinelli2003Geant4,Allison2016Geant4} Monte Carlo
simulation --- their energy, position, direction and particle type --- and then
samples new, physically consistent events at a small fraction of the cost of
running full Geant4 transport.

The package is installed and imported as \code{magi}.

\subsection{What it is for}

The intended use is as a \textbf{drop-in replacement for a Geant4 particle
source}. Run the expensive full-physics simulation once to obtain a reference
sample of particles crossing some surface; train MAGI on that sample; then use
MAGI's output as the source for downstream simulations that would otherwise
require repeating the expensive upstream physics.

This pays off when

\begin{itemize}[nosep]
  \item large numbers of particles are needed downstream;
  \item the events reaching the surface of interest are rare relative to the
    number of primaries thrown, so direct simulation is statistics-starved;
  \item the same downstream simulation is run many times with different
    geometry or configuration choices, and re-deriving the upstream flux each
    time is wasteful.
\end{itemize}

\subsection{What it is not for}

\begin{itemize}[nosep]
  \item Estimating the intensity of an individual spectral line
    (Section~\ref{sec:envelope}).
  \item Detector geometries other than a sphere. The crossing surface is
    hard-coded as a sphere; see Section~\ref{sec:limits}.
  \item Extrapolating outside the energy range or particle mix of the training
    sample. MAGI is a density model of the data it was shown, not a physics
    engine.
\end{itemize}

\subsection{How to read this manual}

Section~\ref{sec:structure} describes the code and what each module is for.
Section~\ref{sec:usage} is the API reference: functions, their settings and
the contracts between them. Section~\ref{sec:tools} covers the command-line
tools around a run. Sections~\ref{sec:fullrun} and~\ref{sec:validation} are
worked examples of a full run and of a validation.
Section~\ref{sec:envelope} states what has actually been measured, and
Section~\ref{sec:limits} lists the traps.

If you are new, start with \code{Example\_Usage.ipynb} at the repository root:
it is a runnable, commented walkthrough of the whole pipeline on synthetic
data, and it finishes in a few minutes.

\begin{maginote}[Conventions]
Energies are in MeV and positions in mm throughout, matching Geant4's own
output convention. Code that must run on the CPU is shown with the
\code{CUDA\_VISIBLE\_DEVICES=-1} prefix, which must be set \emph{before}
\code{magi} is imported.
\end{maginote}

% =====================================================================
\section{Theoretical foundation}
\label{sec:theory}

This section states what MAGI is built on and where each piece is implemented.
It is background, not a derivation; the primary sources are cited so a reader
can go to them directly.

\begin{magiwarn}[Bibliography verification status]
The entries in \code{docs/manual/magi\_theory.bib} are marked
\code{\% VERIFY}: their arXiv identifiers and titles are believed correct, but
no author has yet opened each source to confirm author lists, venue, volume and
pages. Confirm them on NASA ADS, INSPIRE or DBLP before any of this text is
reused in a citable document.
\end{magiwarn}

\subsection{Conditional variational autoencoders}

A variational autoencoder~\cite{Kingma2014VAE} models a density $p(x)$ through
a latent variable $z$, training an encoder $q(z|x)$ and decoder $p(x|z)$
against the evidence lower bound

\begin{equation}
\mathcal{L} \;=\; \mathbb{E}_{q(z|x)}\!\left[\log p(x \mid z)\right]
   \;-\; \mathrm{KL}\!\left(q(z \mid x)\,\|\,p(z)\right).
\end{equation}

The conditional variant~\cite{Sohn2015CVAE} conditions every term on an
observed covariate $c$, giving $q(z|x,c)$, $p(x|z,c)$ and $p(z|c)$. In MAGI $c$
is the particle-type one-hot, which is what lets generation be steered to a
chosen particle mix rather than only reproducing the training proportions.

Implemented in \code{magi/core/model.py}; the reconstruction and KL terms in
\code{magi/core/losses.py}.

\subsection{Normalizing flows for the energy continuum}

A normalizing flow~\cite{Papamakarios2021Flows,Kobyzev2021Flows} represents a
density as an invertible, differentiable map $f$ applied to a simple base
density, with the change-of-variables formula supplying an exact likelihood:

\begin{equation}
\log p(y) \;=\; \log p_u\!\left(f(y)\right)
   \;+\; \log\left|\det \frac{\partial f}{\partial y}\right|.
\end{equation}

MAGI uses monotonic rational-quadratic spline
transforms~\cite{Durkan2019NeuralSplineFlows} for the energy continuum. The
choice matters here for a specific reason: the crossing spectrum of a cosmic-ray
source spans roughly nine decades and contains sharp non-Gaussian structure --- a
Compton edge, a muon peak, a steep low-energy tail --- that a single Gaussian
cannot represent at all, and that an affine flow represents only poorly. A
spline flow is flexible enough to follow that structure while remaining exactly
invertible, which the exact-likelihood training requires.

Implemented in \code{magi/core/flows.py} (\code{ConditionalRQSFlow}), selected
with \code{continuum\_mode="flow"}.

\subsection{The gated line--continuum mixture}

The central modelling decision in v0.8 is that spectral line positions are
\emph{not learned}. A categorical energy head can never place a line more
precisely than one bin, and on a log grid spanning nine decades a bin is
hundreds of eV --- orders of magnitude wider than a real fluorescence line.
Positions therefore become fixed physical inputs, measured from the data and
taken from evaluated atomic data, and the head becomes a mixture:

\begin{equation}
p(y \mid z, c) \;=\; g_0(z,c)\, p_{\text{cont}}(y \mid z, c)
  \;+\; \sum_{\ell=1}^{L} g_\ell(z,c)\,
        \mathcal{N}\!\left(y \,;\, y_\ell,\, \sigma_\ell^2\right),
\qquad \sum_{\ell=0}^{L} g_\ell = 1,
\end{equation}

with $y = \log_{10}(E/\text{MeV})$. The widths $\sigma_\ell$ are pinned to the
instrument's energy resolution rather than fitted.

The gate $g$ faces a severe class imbalance: a line can be $10^{-5}$ of the
sample, so an unweighted cross-entropy is minimized well by never routing
anything to it. MAGI therefore applies focal down-weighting of easy
examples~\cite{Lin2017FocalLoss}, exposed as \code{gate\_focal\_gamma}.

\begin{maginote}[An empirical caution on $\gamma$]
Use $\gamma = 1$. At $\gamma \geq 2$ the gate over-corrects and begins routing
genuine continuum events into the line components, visibly digging a hole in
the continuum on either side of each line. This is measured, not assumed: the
near-line continuum ratio degrades monotonically across a $\gamma$ sweep.
\end{maginote}

\subsection{The learned conditional prior}

With a standard $p(z) = \mathcal{N}(0,I)$, two problems arise. First, the
decoder alone must carry every correlation between energy and geometry.
Second, and more subtly, the model is trained on $z \sim q(z|x)$ but generates
from $z \sim p(z)$; when the aggregated posterior does not match the prior, the
decoder is evaluated at generation time in regions of latent space it was never
trained on~\cite{Tomczak2018VampPrior}. On the reference sources this gap is
large --- validation KL of 21--28 nats at $\texttt{latent\_dim}=8$.

MAGI replaces the fixed prior with a learned conditional coupling flow
$p(z \mid c)$, built from affine coupling layers~\cite{Dinh2017RealNVP} and
trained with a Monte-Carlo estimate of the KL term. This is what brings the
coupling residuals of Section~\ref{sec:envelope} inside their bar, and it is
the difference the v0.7.2 categorical head cannot make up.

In v0.8.2 the prior's conditioning is widened from the particle-type one-hot
alone to $[\,c,\,\text{zone}\,]$, where zone is the continuum/line routing
target (\code{prior\_zone\_conditioning=True}). Without the zone in the
conditioning, the prior has no way to express which mixture component an event
belongs to, so at generation time rare lines are routed by whatever the type
marginal happens to imply.

Implemented in \code{magi/core/priors.py} (\code{ConditionalCouplingPrior}).

\subsection{Atomic transition energies}

Line positions are only as good as the table they are pinned at. MAGI takes
transition \emph{labels} (K$\alpha_1$, K$\beta$, and so on) from
Bearden's compilation~\cite{Bearden1967XRayWavelengths}, but transition
\emph{energies} from the EADL evaluated library~\cite{Perkins1991EADL}, because
EADL is what the Geant4 physics list actually emitted.

This distinction is not pedantic. Pinning to Bearden energies while the
simulation emits EADL displaces Cu~K$\alpha_1$ by 42~eV and Cu~K$\beta$ by
40~eV --- around ten times the 4~eV instrument FWHM the generated lines are
given. The generated and real lines then have essentially no overlap, and no
amount of gate tuning can recover it. Every pinned position is audited against
a centroid measured from the real spectrum
(\code{tools/line\_centroid\_audit.py}).

\subsection{Optimization}

Training uses Adam~\cite{Kingma2015Adam} with a task-adaptive loss-weight
scheduler (\code{magi/training/adaptive\_callbacks.py}) that rebalances the
per-variable loss terms during training, so that no single head dominates the
objective purely because its natural scale is larger.

\subsection{Relation to other work}

MAGI's problem --- learn a compact, resampleable particle source on a surface so
that a multi-stage Monte Carlo can be decoupled --- is the same one
KDSource~\cite{Schmidt2022KDSource} addresses with kernel density estimation.
MAGI substitutes a deep conditional generative model, which scales better to
the joint energy--position--direction--type density and conditions naturally on
particle type.

The broader use of deep generative models to replace expensive detector
simulation is well established in high-energy
physics~\cite{Hashemi2024Review}, with GAN-~\cite{Paganini2018CaloGAN} and
flow-based~\cite{Krause2021CaloFlow} calorimeter shower models now in
production use~\cite{ATLAS2024PhotonShower}. MAGI differs in target: rather
than showers inside a calorimeter, it models the particle flux crossing a
surface, for the cryogenic X-ray instrument background
problem~\cite{Lotti2021XIFUBackground}.

% =====================================================================
\section{Structure of the code}
\label{sec:structure}

\subsection{Repository layout}

The repository has two parts: the installable library, and the experiment
material that uses it.

\begin{center}
\begin{tabular}{@{}p{0.32\textwidth}p{0.60\textwidth}@{}}
\toprule
\textbf{Path} & \textbf{Purpose} \\
\midrule
\code{MAGI\_package/} & The installable Python library (\code{import magi}).
  Everything in Section~\ref{sec:usage} lives here. \\
\code{Example\_Usage.ipynb} & Runnable walkthrough on synthetic data. Start
  here. \\
\code{MAGI\_v0\_*.ipynb} & The experiment notebooks. Highest version number is
  the current pipeline; \code{OldNotebooks/} holds deprecated ones. \\
\code{tools/} & Command-line scripts \emph{around} a run rather than inside
  one: full training runs, the acceptance harness, audits, synthetic stress
  tests (Section~\ref{sec:tools}). \\
\code{scripts/generate\_geant\_source.py} & Produces a Geant4-ready source
  file from a trained run, outside a notebook. This is what the companion
  Geant4 project's \code{/generator/mlScript} macro invokes. \\
\code{CandidateLines/} & Candidate spectral-line tables built from a GDML mass
  model, in the JSON form \code{magi.load\_candidate\_energy\_lines} reads. \\
\code{TrainingData/} & Cleaned, whitespace-delimited crossing data. Gitignored
  --- large binaries stay local. \\
\code{trained\_models/}, \code{checkpoints/} & Trained-run artifacts. Weights
  are gitignored; the JSON config/metadata is versioned. \\
\code{docs/} & Development notes, measurement write-ups and plans for the v0.8
  line, plus this manual under \code{docs/manual/}. \\
\bottomrule
\end{tabular}
\end{center}

\subsection{Package layout}

\begin{lstlisting}[language=shell]
magi/
  __init__.py       public API surface (the canonical list of exports)
  magi.py           high-level convenience API: setup, build_model, train_model
  config.py         environment/seed initialization
  core/
    model.py        the CVAE classes, one per version in the lineage
    losses.py       shared NLL / angular / mixture losses
    geometry.py     physical coordinate transforms
    flows.py        conditional rational-quadratic-spline flow (v0.8 continuum)
    priors.py       conditional coupling prior p(z|cond)   (v0.8)
  data/
    io.py           load/save raw detector tables, normalization, line tables
    preprocessing.py physical features, energy binning, line detection,
                     quantile transforms, gate targets
    dataset.py      split, scale, one-hot conditioning, tf.data construction
  training/
    train.py            compile/fit wrappers
    adaptive_callbacks.py task-adaptive loss rebalancing and monitors
    checkpointing.py    save/load weights + metadata + transformers as a unit
  generation/
    sampling.py       draw latent codes and conditioning
    reconstruction.py invert model outputs back to physical quantities
    export.py         write Geant4 source files (text or chunked binary)
  validation/
    metrics.py      Wasserstein, line-integral recovery
    compare.py      histogram-residual comparisons, range/norm checks
  utils/
    plotting.py         all plot_* helpers, light/dark theme
    model_inspection.py introspect a built model
    circuit_viz.py       interactive routing-circuit HTML (Section~\ref{sec:visualization})
    full_circuit.py      interactive per-event gradient-circuit HTML (Section~\ref{sec:visualization})
tests/              135 tests; see Section~\ref{sec:tools}
docs/USAGE.md       short-form user guide
\end{lstlisting}

\subsection{The data flow}

Every run follows the same seven stages. Each arrow is a function whose output
is the next one's input; each stage has a matching \code{report\_*} or
\code{print\_*} helper that prints sanity checks and returns nothing.

\begin{enumerate}[nosep]
  \item \textbf{Load} --- \code{load\_detector\_table} reads the raw crossing
    file into a DataFrame.
  \item \textbf{Physical features} --- \code{build\_physical\_features} turns
    Cartesian position/direction into the sphere-surface parametrization
    $(u_r, \phi_r, u_v, \phi_v)$ and measures the primary fraction.
  \item \textbf{Feature frame} --- \code{build\_feature\_dataframe} bins
    energy, detects spectral lines and fits the quantile geometry transforms.
  \item \textbf{Dataset} --- \code{filter\_particle\_types\_continuous\_geometry}
    $\rightarrow$ \code{split\_feature\_data} $\rightarrow$
    \code{scale\_continuous\_features} $\rightarrow$
    \code{build\_conditioning\_and\_weights} $\rightarrow$
    \code{build\_tf\_datasets}.
  \item \textbf{Train} --- a class from \code{magi.core}, then
    \code{compile\_model} and \code{fit\_model}.
  \item \textbf{Checkpoint} --- \code{save\_final\_trained\_model} writes
    weights, config, metadata, history, task weights and a summary as one unit.
  \item \textbf{Generate and validate} --- \code{generate\_latent\_outputs}
    $\rightarrow$ \code{reconstruct\_generated\_features} $\rightarrow$
    \code{reconstruct\_generated\_physics} $\rightarrow$ export, with
    \code{magi.validation} scoring the result.
\end{enumerate}

\subsection{Why per-variable heads}

MAGI does not predict one flat output vector. Each physical variable gets its
own head, because each has a different pathology:

\begin{itemize}[nosep]
  \item $u_v$ is bounded on $[-1,1]$ and piles up at the edges. A Gaussian head
    fits it badly, so the model works in $s = \operatorname{atanh}(u)$ space,
    or on a quantile transform of it.
  \item $\phi_r$, $\phi_v$ are periodic. The v0.7 heads predict a
    $(\cos,\sin)$ pair softly regularized onto the unit circle; v0.7.2 predicts
    the angle through a quantile transform directly.
  \item Energy spans many decades and contains narrow spectral lines. This is
    the axis the version lineage is really about
    (Section~\ref{sec:modelvariants}).
\end{itemize}

\subsection{Model variants}
\label{sec:modelvariants}

The classes in \code{magi.core} form a lineage. A trained run records which one
it is in \code{model\_config["model\_class"]}; check that before assuming
behaviour.

\begin{center}
\begin{tabular}{@{}p{0.40\textwidth}p{0.12\textwidth}p{0.42\textwidth}@{}}
\toprule
\textbf{Class} & \textbf{Version} & \textbf{Energy / geometry} \\
\midrule
\code{CVAE\_CatEnergy\_CatUV} & v0.6 &
  Categorical energy, discretized $u_v$. \\
\code{CVAE\_CatEnergy\_CatUV\_TaskAdaptive} & v0.6 &
  As above, plus task-adaptive loss weighting. \\
\code{CVAE\_CatEnergy\_ContGeom\_TaskAdaptive} & v0.7 &
  Continuous geometry, $(\cos,\sin)$ angles. \\
\code{CVAE\_CatEnergy\_ContPhi\_TaskAdaptive} & v0.7.2 &
  Continuous $\phi_r,\phi_v$. The stable fallback head. \\
\code{CVAE\_MixEnergy\_ContPhi\_TaskAdaptive} & v0.8 &
  v0.7.2 geometry; energy becomes a gated mixture of a continuum density and
  fixed-position line components. \\
\bottomrule
\end{tabular}
\end{center}

\paragraph{The v0.8 energy head.} A categorical head can never place a spectral
line more precisely than one bin, and on a log grid spanning nine decades a bin
is hundreds of eV wide --- far broader than a real line. v0.8 therefore stops
asking the model to learn line positions at all. They become \emph{fixed
physical inputs}, measured from the data, and the head becomes a mixture:

\begin{equation}
p(y \mid z, c) \;=\; g_0(z,c)\, p_{\text{cont}}(y \mid z, c)
  \;+\; \sum_{\ell=1}^{L} g_\ell(z,c)\,
        \mathcal{N}\!\left(y \,;\, y_\ell,\, \sigma_\ell^2\right),
\end{equation}

where $y = \log_{10}(E/\text{MeV})$, the gate $g$ is a softmax over
$L+1$ components, the line positions $y_\ell$ are pinned at measured energies
and the widths $\sigma_\ell$ are pinned to the instrument resolution. The
continuum $p_{\text{cont}}$ is either a single Gaussian or a conditional
rational-quadratic-spline normalizing flow (\code{continuum\_mode="flow"}).

\paragraph{The learned prior.} With a standard $\mathcal{N}(0,I)$ prior the
decoder must carry every correlation between energy and geometry on its own.
Setting \code{prior="coupling"} replaces it with a learned conditional coupling
flow $p(z \mid c)$, trained with a Monte-Carlo KL estimate. This is what makes
the coupling residuals in Section~\ref{sec:envelope} pass, and it is v0.8's
main scientific contribution.

Use \code{magi.print\_model\_structure(model)} to print a built model's
generative structure, formulas, configured line table and parameter counts.

% =====================================================================
\section{Usage: functions and settings}
\label{sec:usage}

Every exported function carries a full NumPy-style docstring, so hovering over
a name in VSCode shows its parameters and return value. This section covers the
contracts \emph{between} functions, which a per-function docstring cannot.

\subsection{Installation}

\begin{lstlisting}[language=shell]
pip install -e MAGI_package/      # develop the library
pip install -r requirements.txt   # or install it as a dependency
\end{lstlisting}

Core dependencies: \code{tensorflow} (plus \code{tensorflow-metal} on Apple
Silicon), \code{tensorflow-probability}, \code{scikit-learn}, \code{joblib},
\code{astropy}, \code{optuna}.

\subsection{Environment setup}

\begin{lstlisting}[language=Python]
import magi
magi.initialize_environment(seed=42, cpu_only=True, quiet=True)
magi.print_tf_info()
\end{lstlisting}

\begin{magiwarn}[cpu\_only must come first]
\code{cpu\_only=True} only takes effect if it runs before TensorFlow
initializes its devices. In a script that imports \code{magi} at module level
this is already too late --- set the environment variable instead, before any
import:

\begin{lstlisting}[language=Python]
import os
os.environ["CUDA_VISIBLE_DEVICES"] = "-1"
import magi
\end{lstlisting}

Every script in \code{tools/} does exactly this, and for a reason: on this
workload \code{tensorflow-metal} is slower than the CPU, and importing
\code{tensorflow\_probability} initializes the Metal device, after which
\code{set\_visible\_devices} silently stops working.
\end{magiwarn}

\subsection{Loading data}

\begin{lstlisting}[language=Python]
df = magi.load_detector_table("TrainingData/alloutputDSCryoSphereCR.dat")
magi.report_basic_table_checks(df)
\end{lstlisting}

Three column schemas are auto-detected from the field count on the first data
line:

\begin{center}
\begin{tabular}{@{}clp{0.52\textwidth}@{}}
\toprule
\textbf{Cols} & \textbf{Schema} & \textbf{Adds} \\
\midrule
9 & legacy & \code{EventId ParticleName Energy X Y Z Vx Vy Vz} \\
10 & PrimBool-tagged & \code{PrimBool} (Geant4 \code{ParentID==0}) \\
13 & lineage & \code{ParticleId}, \code{ParentParticleId},
  \code{CreatorProcessName} --- required for radioactive sources \\
\bottomrule
\end{tabular}
\end{center}

\subsection{Preprocessing}

\begin{lstlisting}[language=Python]
prep = magi.build_physical_features(
    df, center=(0.0, 0.0, -507.66), radius=100.0,
    source_type="cosmic_ray",
)
magi.print_physical_summary(prep)

feature_pack = magi.build_feature_dataframe(
    prep,
    energy_binning_mode="log_fixed_count",
    n_bins=512,
    min_counts=20,
    geometry_transform="quantile_u_r_u_v_phi_r_phi_v",
    energy_transform="log10",     # required by the v0.8 mixture head
)
magi.report_feature_dataframe(feature_pack)
\end{lstlisting}

\paragraph{\code{source\_type} decides flux normalization.} With
\code{"cosmic\_ray"} a crossing counts as primary iff \code{PrimBool == 1}.
With \code{"radioactive"} it counts as primary iff
\code{CreatorProcessName == "RadioactiveDecay"}, because for a source embedded
in bulk material the literal Geant4 primary is the decaying nucleus, which
never propagates --- \code{PrimBool} is always~0. Choosing wrongly does not
error; it silently yields a wrong normalization factor.

\paragraph{\code{geometry\_transform} is a stringly-typed contract.} The value
(\code{"arctanh\_uv\_discrete"}, \code{"quantile\_u\_r\_u\_v"},
\code{"quantile\_u\_r\_u\_v\_phi\_r\_phi\_v"}) fixes the column ordering shared
by preprocessing, the model's input layout and reconstruction. It is persisted
in the saved metadata and \textbf{must} match between the run that trained a
model and the run that generates from it. A mismatch does not raise --- it
produces physically wrong output.

\subsection{Spectral lines}

Line handling is the part of the pipeline most specific to v0.8, and the part
where a silent mistake is most expensive.

Note the ordering: lines are detected on the \emph{raw} energies, on their own
fine binning, \emph{before} \code{build\_feature\_dataframe} is called. The
feature frame's binning is a separate, coarser thing --- with the mixture head
it no longer constrains where a line can be placed, only how the continuum is
described. The resulting gate targets are then carried into the dataset as
extra feature columns.

\begin{lstlisting}[language=Python]
payload = magi.load_candidate_energy_lines("CandidateLines/cryosphere_lines.json")
candidate_lines = payload["lines"]

res = magi.detect_energy_lines(
    E, binning_mode="log_fixed_count", n_bins=1024,
    candidate_lines=candidate_lines,
    refine_bin_width_mev=4.0e-6)   # refine to the pinned line width
magi.print_detected_energy_lines(res)

matched = [m for m in res["matched_lines"] if m["count"] >= 100]
matched += magi.confirm_unresolved_candidate_lines(
    E, candidate_lines, matched, resolution_ev=4.0)

gate_targets = magi.build_gate_targets(
    E, feature_pack["energy_bins"], matched,
    bandwidth_mode="resolution",       # the default; see the warning below
    bandwidth_fwhm_mev=4.0e-6)
line_logsigma = magi.line_logsigma_from_resolution(
    [m["candidate_energy_mev"] for m in matched], resolution_ev=4.0, fwhm=True)
\end{lstlisting}

\begin{magiwarn}[\code{bandwidth\_mode}: a case where the more correct label trained worse]
\code{bandwidth\_mode} controls how \code{build\_gate\_targets} decides which
\emph{real} events count as line events when building the gate's training
target. It does not affect the generated line width, which is pinned separately
via \code{line\_logsigma\_from\_resolution}.

Raw Geant4 crossing data has \textbf{no detector response applied}, so simulated
fluorescence lines are exactly monoenergetic: in the reference cosmic-ray
source all 7\,542 Cu~K$\alpha_1$ events share one identical \code{float64}
energy, and every nearby continuum event is a singleton. On that reasoning,
using the instrument's 4~eV FWHM as a Gaussian kernel to decide line membership
is the wrong tool --- it hands weight to continuum a few eV away and blurs
doublets closer than a few resolution widths --- and
\code{bandwidth\_mode="exact"} was added to match the simulation's own ground
truth.

\textbf{Measured on the cosmic-ray source, it made things worse.} It did not
fix the Cu~K$\alpha_1$ overshoot it was built for (2.011 versus 2.052, inside
the noise), it regressed the previously passing Al~K$\alpha_1$ from
$0.959 \pm 0.025$ to 0.578, it regressed Cu~K$\beta$ from $1.034 \pm 0.050$ to
0.753, and it pushed the coupling residual from $0.0285 \pm 0.0101$ to 0.054,
outside its bar. The tight tolerance also drove Cu~K$\alpha_2$'s zone
probability to zero, effectively deleting it as a modelled component.

\code{"resolution"} is therefore the default and is what every confirmed result
in Section~\ref{sec:envelope} was produced with. \code{"exact"} remains
available and tested, but it is a research option, not a recommendation.

The lesson is worth stating plainly: a training \emph{label} that is more
faithful to the simulation is not automatically a label that trains better. The
resolution kernel appears to be doing useful work as a soft target, and
replacing it with a hard one removed that.
\end{magiwarn}

\code{confirm\_unresolved\_candidate\_lines} exists because a doublet closer
together than one detection bin is always a single peak, so only one member is
ever detected regardless of how good the matching logic is --- the gap is in
peak detection, not matching. It re-measures every unmatched candidate at eV
resolution and confirms the ones that are real. Both
\code{tools/run\_v0\_8\_real.py} and \code{tools/acceptance\_v0\_8.py} must
call it identically, or the checkpoint's \code{n\_lines} silently misaligns
against the rebuilt pipeline's line count.

\subsection{Building the dataset}

\begin{lstlisting}[language=Python]
dataset_pack   = magi.filter_particle_types_continuous_geometry(
                     feature_pack["feat"], prob_threshold=1e-5)
split_pack     = magi.split_feature_data(dataset_pack, random_state=42)
scaled_pack    = magi.scale_continuous_features(split_pack, scale_cols=())
condition_pack = magi.build_conditioning_and_weights(
                     scaled_pack, dataset_pack["n_types"], alpha=0.5)
ds_pack        = magi.build_tf_datasets(condition_pack, batch_size=4096)
magi.report_tf_datasets(ds_pack)
\end{lstlisting}

Note \code{scale\_cols=()} for v0.7+ quantile geometry: those columns are
already standardized by their own transform, and scaling them twice is a
mistake the empty tuple prevents. The split is stratified on particle type;
\code{alpha} controls how strongly rare species are up-weighted
($0$~=~natural proportions, $1$~=~fully balanced, $0.5$~=~default).

\subsection{Training}

\begin{lstlisting}[language=Python]
model = magi.CVAE_MixEnergy_ContPhi_TaskAdaptive(
    n_cont=4, n_types=dataset_pack["n_types"], latent_dim=8,
    line_positions_y=line_positions_y,
    line_logsigma_init=line_logsigma,
    continuum_mode="flow",
    continuum_flow_warp="cdf",
    energy_flow_condition="z_cond",
    prior="coupling",
    prior_zone_conditioning=True,
    zone_probs=zone_probs,
    gate_focal_gamma=1.0,
    w_gate_aux=2.0,
)

magi.compile_model(model, learning_rate=2e-4, optimizer="adam")
history = magi.fit_model(
    model, ds_pack["train_ds"], ds_pack["val_ds"],
    epochs=40, callbacks=magi.build_default_callbacks(), verbose=2)
\end{lstlisting}

\begin{center}
\begin{tabular}{@{}p{0.30\textwidth}p{0.14\textwidth}p{0.48\textwidth}@{}}
\toprule
\textbf{Setting} & \textbf{Default} & \textbf{Effect} \\
\midrule
\code{latent\_dim} & 8 & Latent width. Larger widens the
  posterior/prior gap that the coupling prior exists to close. \\
\code{continuum\_mode} & \code{"gaussian"} & Use \code{"flow"}: a single
  Gaussian cannot represent a multi-decade continuum. \\
\code{continuum\_flow\_warp} & \code{"affine"} & Use \code{"cdf"}: pre-warps
  the flow's input by the empirical CDF, which is what recovers the Compton
  edge and the low-energy tail. \\
\code{energy\_flow\_condition} & --- & \code{"z\_cond"} conditions the
  continuum flow on both latent and conditioning. \\
\code{prior} & \code{"gaussian"} & Use \code{"coupling"} for the learned
  conditional prior. This is what makes coupling residuals pass. \\
\code{prior\_zone\_conditioning} & \code{False} & Condition the prior on
  \code{[type, zone]} rather than type alone. Requires \code{zone\_probs}. \\
\code{gate\_focal\_gamma} & 0 & Focal weighting on the gate loss.
  Use~1; $\gamma \geq 2$ digs the continuum out from between the lines. \\
\code{w\_gate\_aux} & --- & Weight of the auxiliary gate-routing loss. \\
\bottomrule
\end{tabular}
\end{center}

\subsection{Checkpointing}

\begin{lstlisting}[language=Python]
magi.save_final_trained_model(
    model=model,
    save_dir="trained_models/my_run",
    model_name="mix_CR",
    history=history,
    model_config=model.to_generation_config(),
    preprocessing_metadata=preprocessing_metadata,
)
joblib.dump(transformers, "trained_models/my_run/mix_CR_quantile_transformers.joblib")
\end{lstlisting}

\begin{magiwarn}[A checkpoint is one unit]
A run directory holds \code{*.weights.h5}, \code{*\_config.json},
\code{*\_metadata.json}, \code{*\_history.json}, \code{*\_task\_weights.json},
\code{*\_summary.txt} and --- for quantile-geometry runs ---
\code{*\_quantile\_transformers.joblib}. Generation needs the config and
metadata as much as the weights, because they carry the transforms that invert
the preprocessing. Copy or move them together.
\end{magiwarn}

Always pass \code{model.to\_generation\_config()} rather than a hand-written
dict. For v0.8 heads the loader \emph{requires} every key it emits and raises
on a missing one. This guard exists because the old behaviour was
\code{model\_config.get(key, default)} everywhere, and the defaults
(\code{continuum\_mode="gaussian"}, \code{prior="gaussian"},
\code{continuum\_flow\_warp="affine"}) select a \emph{code path} rather than a
layer shape --- so a renamed key produced a silently different architecture
that loaded without complaint and generated wrong output.

Checkpoints from v0.8.2 carry \code{config\_version: 2}. Older
\code{config\_version: 1} checkpoints load unchanged:
\code{prior\_zone\_conditioning} defaults to \code{False}, reconstructing the
exact pre-existing architecture.

\subsection{Generation}

\begin{lstlisting}[language=Python]
model = magi.load_task_adaptive_model_for_generation(
    save_dir="trained_models/my_run", model_name="mix_CR",
    model_config=magi.load_json("trained_models/my_run/mix_CR_config.json"),
    n_types=n_types, type_weights=type_weights, radius=100.0)

gen_pack = magi.generate_latent_outputs(
    model, n_samples=n_real, type_probs=type_probs,
    n_types=n_types, idx_to_type=idx_to_type)

reco = magi.reconstruct_generated_features(
    gen_pack, energy_head_mode="mixture",
    geometry_metadata=geometry_metadata, energy_metadata=energy_metadata)

gen = magi.reconstruct_generated_physics(
    reco, center=(0.0, 0.0, -507.66), radius=100.0)
\end{lstlisting}

\begin{magiwarn}[Generate as many events as you have real ones]
The line-recovery metrics compare raw counts. Generating fewer events than the
real sample deflates every ratio by exactly the shortfall --- a $3.44\times$
error on the reference cosmic-ray source. Set \code{n\_samples} to the real
event count when validating.
\end{magiwarn}

For a Geant4-ready file, use the one-call entry point, which loads the
checkpoint and streams events to disk in chunks:

\begin{lstlisting}[language=Python]
magi.generate_detector_input_file(
    save_dir="trained_models/my_run", model_name="mix_CR",
    model_config=model_config, n_types=n_types,
    type_weights=type_weights, type_probs=type_probs,
    idx_to_type=idx_to_type, geometry_metadata=geometry_metadata,
    energy_head_mode="mixture", energy_metadata=energy_metadata,
    output_file="generated/source.dat", n_events=10_000_000,
    radius=100.0, center=(0.0, 0.0, -507.66),
    output_format="binary", seed=42)
\end{lstlisting}

Neutrinos are filtered out by default (they do not contribute to detector
background here). With \code{generate\_until\_n\_written=True}, the default,
generation continues until \code{n\_events} \emph{transportable} particles have
been written, so Geant4 receives exactly the requested count.

\subsection{Visualization}
\label{sec:visualization}

Two self-contained, interactive HTML tools in \code{magi.utils} inspect a
trained v0.8 mixture run without a plotting library or a running kernel ---
each call writes one \code{.html} file that opens in any browser, with its own
light/dark palette and no external dependency.

\paragraph{\code{save\_routing\_circuit}} The gate's conditional routing, per
type: for each particle type, a fan-out diagram of one-hot conditioning
$\to$ gate $\to$ zone, colored by the checkpoint's own \code{zone\_probs} ---
the same per-type empirical routing frequency the learned conditional prior
is conditioned on (\S\ref{sec:theory}).

\begin{lstlisting}[language=Python]
magi.save_routing_circuit(save_dir="trained_models/v0_8_2_priorzone_CR",
                           model_name="mix_CR")
\end{lstlisting}

Reads only the checkpoint's own config/metadata --- no raw data, no loaded
model, so it runs in under a second. Requires
\code{prior\_zone\_conditioning=True} at training time: a \code{config\_version
1} checkpoint (or any run with zone-conditioning off) has no \code{zone\_probs}
to plot and raises \code{KeyError}.

\paragraph{\code{save\_full\_circuit}} A deeper, per-event trace: for one real
held-out event per type, gradient$\times$activation attribution through
\emph{every} stage --- encoder, latent $z$, decoder stem, deep trunk, and all
three heads (energy, position, direction) --- ending in the real marginal
energy spectrum, with the selected type's own contribution shaded on top of
the combined one. The attribution target is the model's own per-event
reconstruction loss, not just the gate logit, so branches with no gradient
path to the gate at all --- the energy branch reads the decoder stem directly,
bypassing the deep trunk that feeds position and direction --- still get
genuine, nonzero attribution. An adjustable ``highlight top $N$\,\%'' slider
and a light/dark toggle are built into the page.

\begin{lstlisting}[language=Python]
magi.save_full_circuit(
    save_dir="trained_models/v0_8_2_priorzone_CR", model_name="mix_CR",
    training_data_filepath="TrainingData/alloutputDSCryoSphereCR.dat",
    candidate_lines_file="CandidateLines/CANDIDATE_ENERGY_LINES_"
                          "SRON_CCNwithXFDM_NoShield_FlowerCryoAC_fixed_EADL.json",
    center=(0.0, 0.0, -507.66), radius=100.0)
\end{lstlisting}

\begin{magiwarn}[This one re-derives your preprocessing --- budget the time]
Unlike the routing circuit, \code{save\_full\_circuit} needs a real held-out
sample to trace, which a checkpoint alone does not carry. It re-runs line
matching, gate-target construction, and the quantile/log10 transforms against
the raw detector table you point it at, then rebuilds the same train/test
split (\code{random\_state=42} by default, matching
\code{tools/run\_v0\_8\_real.py}) before loading the model. Expect roughly one
preprocessing pass over your source: a couple of minutes for the 3.4\,M-event
CR file, longer for a bigger one (Torio's 8.8\,M events takes several minutes
on CPU).
\end{magiwarn}

The stage layout (how many encoder/trunk/branch columns are drawn, and where
they connect) is derived from the loaded checkpoint's own architecture, not
hardcoded to any one source --- it renders correctly whether the checkpoint
has two particle types and four gate zones (Small/K-40) or six and six (CR).

% =====================================================================
\section{The tools}
\label{sec:tools}

Scripts in \code{tools/} operate around a run rather than inside one. All of
them force CPU execution before importing \code{magi}. Run them from the
repository root.

\subsection{Training and scoring}

\paragraph{\code{run\_v0\_8\_real.py}} A full training run on the real sources
with live per-epoch logging and chunked generation, so a long run is observable
and OOM-safe. Saves a checkpoint, spectra and a recovery report per source.

\begin{lstlisting}[language=shell]
python tools/run_v0_8_real.py --sources CR Small --run-tag v0_8_2 \
    --epochs 40 --seed 42 --prior-zone-conditioning
\end{lstlisting}

Flags: \code{--sources}, \code{--run-tag}, \code{--epochs}, \code{--seed},
\code{--n-gen}, \code{--gen-chunk}, \code{--prior-zone-conditioning},
\code{--gate-target-bandwidth-mode} (default \code{exact}),
\code{--continuum-flow-bins}, \code{--warp-boost-cr}.

\paragraph{\code{acceptance\_v0\_8.py}} The pass/fail harness. It grades a
trained run against explicit thresholds, so an 80-minute run is scored
objectively rather than by eye. The same numbers are the paper's validation
table.

\begin{lstlisting}[language=shell]
python tools/acceptance_v0_8.py --sources CR Small --seeds 42 7 13 \
    --json docs/_data/acceptance.json
\end{lstlisting}

It reports four things: per-line integral recovery (resolution-scaled
$\pm5\sigma$ window, local continuum subtracted, normalized for
$N_{\text{real}}/N_{\text{gen}}$); near-line continuum ratio, i.e. whether the
gate dug a hole beside each line; the energy-marginal Wasserstein distance,
global and per decade band; and the coupling residuals.

\begin{magiwarn}[Never trust a single seed on the cosmic-ray source]
The per-band Wasserstein distance on the CR source has a seed-to-seed
$\sigma/\mu$ of 40--50\,\%, larger than most effects worth chasing. Two changes
were adopted on single-seed evidence during development and both evaporated
under a three-seed grid. Use \code{--seeds 42 7 13} and report mean $\pm$ std.
The Small source is tighter, around 14\,\%.
\end{magiwarn}

\paragraph{\code{score\_v0\_7\_2.py}} Scores a v0.7.2 categorical-head run on
the same axes, so the two heads can be compared like with like. It reads each
checkpoint's own \code{preprocessing\_metadata} rather than re-deriving it.

\subsection{Audits}

\paragraph{\code{line\_centroid\_audit.py}} For each matched line, histograms
the real spectrum locally at eV resolution, fits the centroid and reports
$|\text{catalogue} - \text{measured}|$ in FWHM units, failing loudly above
one~FWHM. Run this before committing to a long training run: a mispositioned
line cannot be fixed by any amount of training.

\paragraph{\code{build\_candidate\_lines\_from\_geant4.py}} Builds the
candidate-line table from a GDML mass model and Geant4's fluorescence data. It
takes transition \emph{labels} from the Bearden table (only that directory
carries them) and transition \emph{energies} from the EADL table, because EADL
is what the simulation actually emitted. Getting this backwards puts every
instrumental line tens of eV off --- up to 10~FWHM.

\paragraph{\code{gate\_posterior\_prior\_audit.py}} Compares the gate's
component fractions under $z \sim q(z|x)$ against $z \sim p(z|\text{cond})$.
A large divergence localizes a line failure to the prior rather than the gate.

\subsection{Synthetic stress tests}

\code{synthetic\_stress\_test\_coupled.py},
\code{\_cr\_multimodal.py} and \code{\_small\_singlepeak.py} build synthetic
data with a known answer --- an injected energy--geometry coupling, a rare-line
cluster, a sparse tail. Use them to gate a change \emph{before} spending a real
run on it; they finish in minutes.

\subsection{Unit tests}

\begin{lstlisting}[language=shell]
python -m pytest MAGI_package/tests/ -q
\end{lstlisting}

135 tests covering flow round-trip identity, checkpoint save/load config
matching, gate-target construction (including the exact/resolution A/B), prior
zone-conditioning, the two visualization tools (\S\ref{sec:visualization}), and
public API docstring coverage. They catch mechanical regressions; they do not
tell you a trained model reproduced your source.

% =====================================================================
\section{Example: a full run}
\label{sec:fullrun}

This is a complete run on the reference cosmic-ray source, from raw data to a
Geant4 source file. On a CPU-only Apple Silicon machine it takes roughly 90
minutes for 40 epochs plus generation.

\subsection{The short path}

Everything below is what \code{tools/run\_v0\_8\_real.py} does. In practice:

\begin{lstlisting}[language=shell]
export CUDA_VISIBLE_DEVICES=-1
python tools/run_v0_8_real.py --sources CR --run-tag my_run \
    --epochs 40 --seed 42 --prior-zone-conditioning \
    2>&1 | tee logs/my_run_CR.log
\end{lstlisting}

\subsection{The long path, step by step}

\begin{lstlisting}[language=Python]
import os
os.environ["CUDA_VISIBLE_DEVICES"] = "-1"
import numpy as np, joblib, magi

magi.initialize_environment(seed=42, cpu_only=True)

# --- 1. load --------------------------------------------------------
df = magi.load_detector_table("TrainingData/alloutputDSCryoSphereCR.dat")
magi.report_basic_table_checks(df)

# --- 2. physical features -------------------------------------------
prep = magi.build_physical_features(
    df, center=(0.0, 0.0, -507.66), radius=100.0, source_type="cosmic_ray")
magi.print_physical_summary(prep)

# --- 3. lines: detect on the RAW energies, before the feature frame --
candidate_lines = magi.load_candidate_energy_lines(
    "CandidateLines/cryosphere_lines.json")["lines"]
E = prep["features"]["Energy"].to_numpy()

res = magi.detect_energy_lines(
    E, binning_mode="log_fixed_count", n_bins=1024,
    prominence_factor=3.0, window=5, candidate_lines=candidate_lines,
    refine_bin_width_mev=4.0e-6)
matched = [m for m in res["matched_lines"] if m["count"] >= 100]
matched += magi.confirm_unresolved_candidate_lines(
    E, candidate_lines, matched, resolution_ev=4.0)
magi.print_detected_energy_lines(res)

# --- 4. feature frame ------------------------------------------------
feature_pack = magi.build_feature_dataframe(
    prep, energy_binning_mode="log_fixed_count", n_bins=512, min_counts=20,
    geometry_transform="quantile_u_r_u_v_phi_r_phi_v",
    energy_transform="log10")          # the v0.8 head needs y = log10(E)
magi.report_feature_dataframe(feature_pack)

E_full = feature_pack["filtered_prep"]["features"]["Energy"].to_numpy()

gate_targets = magi.build_gate_targets(
    E_full, feature_pack["energy_bins"], matched,
    bandwidth_mode="resolution", bandwidth_fwhm_mev=4.0e-6)

line_positions_y = np.log10(
    [m["candidate_energy_mev"] for m in matched]).astype(np.float32)
line_logsigma = magi.line_logsigma_from_resolution(
    10.0 ** line_positions_y, resolution_ev=4.0, fwhm=True)

# --- 5. dataset: the gate targets ride along as extra feature columns -
feat = feature_pack["feat"].copy()
for j in range(gate_targets.shape[1]):
    feat[f"gate_target_{j}"] = gate_targets[:, j]
cont_cols = ("u_r_q", "u_v_q", "phi_r_q", "phi_v_q", "energy_y") + tuple(
    f"gate_target_{j}" for j in range(gate_targets.shape[1]))

dataset_pack   = magi.filter_particle_types_continuous_geometry(
                     feat=feat, prob_threshold=1e-5, cont_cols=cont_cols)
split_pack     = magi.split_feature_data(dataset_pack, random_state=42)
magi.report_split_summary(split_pack, dataset_pack["n_types"])
scaled_pack    = magi.scale_continuous_features(split_pack, scale_cols=())
condition_pack = magi.build_conditioning_and_weights(
                     scaled_pack, dataset_pack["n_types"],
                     idx_to_type=dataset_pack["idx_to_type"], alpha=0.5)
ds_pack        = magi.build_tf_datasets(condition_pack, batch_size=4096)
magi.report_tf_datasets(ds_pack)

# --- 6. per-type zone probabilities for the conditioned prior --------
n_zones   = gate_targets.shape[1]
zone_cols = dataset_pack["X_cont_raw"][:, -n_zones:]
zone_probs = np.zeros((dataset_pack["n_types"], n_zones))
for t in range(dataset_pack["n_types"]):
    mask = dataset_pack["y_type"] == t
    row  = zone_cols[mask].mean(axis=0) if mask.any() else np.zeros(n_zones)
    s    = row.sum()
    zone_probs[t] = row / s if s > 0 else np.eye(n_zones)[0]

# --- 7. model and training -------------------------------------------
model = magi.CVAE_MixEnergy_ContPhi_TaskAdaptive(
    n_cont=4, n_types=dataset_pack["n_types"], latent_dim=8,
    line_positions_y=line_positions_y, line_logsigma_init=line_logsigma,
    continuum_mode="flow", continuum_flow_warp="cdf",
    energy_flow_condition="z_cond",
    prior="coupling", prior_zone_conditioning=True, zone_probs=zone_probs,
    gate_focal_gamma=1.0, w_gate_aux=2.0)

magi.compile_model(model, learning_rate=2e-4)
magi.print_model_structure(model)

history = magi.fit_model(
    model, ds_pack["train_ds"], ds_pack["val_ds"],
    epochs=40, callbacks=magi.build_default_callbacks(), verbose=2)

# --- 8. checkpoint ----------------------------------------------------
magi.save_final_trained_model(
    model=model, save_dir="trained_models/my_run", model_name="mix_CR",
    history=history, model_config=model.to_generation_config(),
    preprocessing_metadata=preprocessing_metadata,
    normalization_metadata=prep["normalization"])
\end{lstlisting}

\subsection{Producing the Geant4 source file}

\begin{lstlisting}[language=shell]
python scripts/generate_geant_source.py \
    --save-dir trained_models/my_run \
    --model-name mix_CR \
    --metadata-file trained_models/my_run/mix_CR_metadata.json \
    --output-file generated/cr_source.dat \
    --n-events 10000000
\end{lstlisting}

The companion Geant4 project consumes this directly through its
\code{/generator/mlScript} macro command.

% =====================================================================
\section{Example: a validation}
\label{sec:validation}

Training finishing is not evidence that it worked. This is how a run gets
graded.

\subsection{The short path}

\begin{lstlisting}[language=shell]
python tools/acceptance_v0_8.py --sources CR --run-tag my_run \
    --seeds 42 7 13 --json docs/_data/my_run_acceptance.json
\end{lstlisting}

\subsection{What it checks, and against what}

\begin{center}
\begin{tabular}{@{}p{0.30\textwidth}p{0.22\textwidth}p{0.42\textwidth}@{}}
\toprule
\textbf{Criterion} & \textbf{Threshold} & \textbf{Notes} \\
\midrule
Per-line recovery & $[0.8,\,1.2]$, two-sided &
  Only for lines above the detection floor. Two-sided because
  over-generation is the dominant failure. \\
Per-line stability & $\sigma \leq 0.15$ over $\geq 3$ seeds &
  So a user's single unattended run likely lands in band. \\
Detection floor & $\geq 5\sigma$ Poisson &
  The particle-physics discovery convention. Below-floor lines are
  \emph{not modelled}, and their absence is a correct result. \\
Near-line continuum & $\geq 0.85$ &
  Did the gate dig a hole beside the line? \\
Energy marginal & Wasserstein $\leq 0.05$ &
  Global, on $\log_{10}E$. \\
Coupling & $\max|\Delta\mathrm{corr}| \leq 0.05$ &
  Over $(\log E, u_r, u_v, \phi_r, \phi_v)$. \\
\bottomrule
\end{tabular}
\end{center}

\subsection{Doing it by hand}

\begin{lstlisting}[language=Python]
real = magi.reconstruct_real_test_physics(
    split_pack["X_cont_test"], split_pack["E_test_raw"],
    center=(0.0, 0.0, -507.66), radius=100.0,
    geometry_metadata=geometry_metadata)

# Generate as many events as there are real ones - see the warning in 3.9.
gen_pack = magi.generate_latent_outputs(
    model, n_samples=real["E"].size, type_probs=type_probs,
    n_types=n_types, idx_to_type=idx_to_type)
reco = magi.reconstruct_generated_features(
    gen_pack, energy_head_mode="mixture",
    geometry_metadata=geometry_metadata, energy_metadata=energy_metadata)
gen  = magi.reconstruct_generated_physics(reco, center=(0,0,-507.66), radius=100.0)

# 1. marginals
scores = magi.compute_wasserstein_scores(real, gen)
print(scores["logE"])

# 2. geometric constraints that must hold by construction
magi.report_generated_constraints(gen, radius=100.0)
magi.report_final_ranges(real, gen)

# 3. per-line recovery
recovery = magi.compute_line_integral_recovery(
    real["E"], gen["E"], matched, energy_bins,
    energy_component_idx_gen=reco["energy_component_idx_gen"],
    resolution_ev=4.0, neighbour_lines=candidate_lines)
for r in recovery:
    print(f"{r['label']:20s} {r['recovery_ratio']:6.3f}  "
          f"sig={r['line_significance']:7.1f}")

# 4. the joint distribution - this is the one that matters
cols = ["logE", "u_r", "u_v", "phi_r", "phi_v"]
df_real, df_gen, df_both = magi.build_real_generated_featureframes(real, gen)
_, corr_real = magi.plot_correlation_matrix(df_real, cols, show=False)
_, corr_gen  = magi.plot_correlation_matrix(df_gen,  cols, show=False)
print("coupling residual:", np.abs(corr_gen - corr_real).to_numpy().max())

magi.compare_hist_with_residuals(real["E"], gen["E"], "Energy [MeV]",
                                 bins=200, ratio_clip=1.0)
\end{lstlisting}

\subsection{Reading the output}

\begin{itemize}[nosep]
  \item \textbf{Coupling residual} is the headline. Below 0.05 means the joint
    distribution is reproduced, which is what MAGI is for.
  \item \textbf{Wasserstein on $\log_{10}E$} below 0.05 means the spectrum's
    shape is right.
  \item \textbf{\code{line\_significance}} below 5 means the line is not
    significantly detected in the real data either; its recovery ratio is
    noise and should be ignored, not reported as a failure.
  \item \textbf{\code{sideband\_contaminated}} true means a neighbouring line
    sits in the continuum side-bands, so the continuum subtraction for that
    line is unreliable.
  \item A statistical floor of
    $\sqrt{1/n_{\text{real}} + 1/n_{\text{gen}}}$ bounds how precisely any
    recovery ratio can be known. For a line with a few hundred events this is
    already several per cent.
\end{itemize}

% =====================================================================
\section{Accuracy: what has actually been measured}
\label{sec:envelope}

All numbers are mean $\pm$ std over three seeds (42/7/13) on the two reference
sources, produced by \code{tools/acceptance\_v0\_8.py}.

\subsection{Validated}

\begin{center}
\begin{tabular}{@{}lccc@{}}
\toprule
\textbf{Quantity} & \textbf{CryoSphere-CR} & \textbf{CryoSphere-Small} & \textbf{Bar} \\
\midrule
Coupling, $\max|\Delta\mathrm{corr}|$ & $0.0285 \pm 0.0101$ & $0.0357 \pm 0.0102$ & $\leq 0.05$ \\
Wasserstein on $\log_{10}E$           & $0.0241 \pm 0.0046$ & $0.0065 \pm 0.0006$ & $\leq 0.05$ \\
\bottomrule
\end{tabular}
\end{center}

The joint distribution holds up with error bars: every real cross-correlation
over $(\log E, u_r, u_v, \phi_r, \phi_v)$ is reproduced inside the bar, on both
sources, across seeds. The v0.7.2 categorical head does not achieve this --- its
coupling residual is 0.226, eight times worse.

\subsection{Not validated}

Per-line \emph{intensities} are wrong by factors of roughly $0.7\times$ to
$4.9\times$, and unstable across seeds.

\begin{center}
\begin{tabular}{@{}llc@{}}
\toprule
\textbf{Line} & \textbf{Recovery} & \textbf{Status} \\
\midrule
CR Al K$\alpha_1$      & $0.959 \pm 0.025$ & in band \\
CR Cu K$\beta$         & $1.034 \pm 0.050$ & in band \\
CR $e^+e^-$ 511 keV    & $1.825 \pm 0.324$ & over-generated, unstable \\
CR Cu K$\alpha_1$      & $2.052 \pm 0.034$ & over-generated \\
Small Al K$\alpha_2$   & $3.412 \pm 0.376$ & far over \\
Small Cu K$\alpha_1$   & $4.903 \pm 0.356$ & far over \\
\bottomrule
\end{tabular}
\end{center}

Practical consequences:

\begin{itemize}[nosep]
  \item Line \emph{positions} are reliable --- pinned physical inputs, audited
    to $\leq 0.5$~eV against the measured spectrum on all sources. It is the
    number of events in each line that is wrong.
  \item Do not use generated output to estimate a fluorescence or annihilation
    line flux, or any quantity dominated by one.
  \item The continuum between lines, and the total energy marginal, are within
    the bar.
  \item Rare lines fare worse than strong ones, and the failure is
    source-dependent. Measure on \emph{your} source with the acceptance
    harness rather than assuming these numbers transfer.
\end{itemize}

\subsection{Choosing between the heads}

Neither head is correct for lines today. v0.7.2 has a better energy marginal on
the cosmic-ray source (0.0141 versus 0.0241) but produces \textbf{zero} events
in the 511~keV window against 52\,225 real, and its coupling is eight times
worse. If your use depends on correlations, use v0.8.2. If it depends only on
the marginal spectrum and you want the most conservative option, v0.7.2 remains
available.

\subsection{Negative results already established}

Recorded so they are not re-attempted without new information:

\begin{itemize}[nosep]
  \item \textbf{Per-line gate class weights}, calibrated from measured routing:
    no effect distinguishable from noise. The untouched control line moved more
    ($+27\,\%$) than either weighted line.
  \item \textbf{Continuum polish} (CDF-warp knot boost, more flow bins):
    a three-seed grid put every difference under $1.2\sigma$ of the noise.
  \item \textbf{Exact-match gate labels}
    (\code{bandwidth\_mode="exact"}): physically the more faithful label, but
    it did not fix the line it was built for and broke a confirmed result and
    the coupling bar. See the warning in Section~\ref{sec:usage} and
    \code{docs/v0.8.1\_line\_truth.md}~\S14.2.
\end{itemize}

% =====================================================================
\section{Limitations and traps}
\label{sec:limits}

\begin{magiwarn}[Geometry is hard-limited to a sphere]
\code{build\_physical\_features}'s \code{center}/\code{radius} and the
transforms in \code{core/geometry.py} all assume the crossing surface is a
sphere. Adapting MAGI to a planar, cylindrical or box surface is \textbf{not} a
matter of different parameters: it requires reworking the geometry-transform
layer in both \code{preprocessing.py} and \code{generation/reconstruction.py},
which must stay in sync.
\end{magiwarn}

\begin{description}[style=unboxed,leftmargin=0pt,itemsep=4pt]
\item[The training data is a crossing spectrum, not deposited energy.]
  Events are recorded where they cross a virtual sphere, with no detector
  response applied. Classical detector escape peaks may legitimately be absent;
  a null result there is a real result.

\item[\code{geometry\_transform} mismatches do not raise.]
  They produce physically wrong output. The string is persisted in metadata for
  this reason --- pass the metadata rather than re-specifying by hand.

\item[Particle-type filtering silently drops rare species.]
  If a rare-but-relevant type is missing from generated output, check
  \code{prob\_threshold} first.

\item[Energy binning limits which lines can be \emph{found}.]
  With the v0.8 mixture head, binning no longer limits where a line can be
  \emph{placed}, but on a log grid spanning many decades a bin can be hundreds
  of eV wide, so close doublets fall in one bin. Use
  \code{refine\_bin\_width\_mev} and
  \code{confirm\_unresolved\_candidate\_lines}, then verify with
  \code{tools/line\_centroid\_audit.py}.

\item[Quantile transforms need enough samples.]
  Small datasets give unreliable tail behaviour in the fitted transform, which
  shows up as implausible extreme values in generated output.

\item[Flux normalization depends on \code{source\_type}.]
  Choosing wrongly does not error. Radioactive sources additionally need the
  13-column lineage schema.

\item[Unit tests cover machinery, not physics.]
  118 passing tests do not tell you a trained model reproduced your source.
  That is what Section~\ref{sec:validation} is for.

\item[CPU-only by default.]
  On this workload \code{tensorflow-metal} is slower than the CPU.
\end{description}

% =====================================================================
\section{Where to look next}

\begin{center}
\begin{tabular}{@{}p{0.42\textwidth}p{0.52\textwidth}@{}}
\toprule
\textbf{Document} & \textbf{Contents} \\
\midrule
\code{MAGI\_package/docs/USAGE.md} & Short-form version of this manual. \\
\code{Example\_Usage.ipynb} & Runnable walkthrough on synthetic data. \\
\code{docs/v0.8.1\_line\_truth.md} & The measurement record behind
  Section~\ref{sec:envelope}: what was tried, what it showed. \\
\code{docs/v0.8.2\_RoadmapForAdoption.md} & The remaining plan and the
  acceptance bar as restated. \\
\code{docs/v0.8\_learnable\_prior\_theory.md} & Why the conditional prior was
  needed. \\
\code{docs/v0.8\_v072\_comparison.md} & Head-to-head history. Note its absolute
  numbers predate the metric fix. \\
\bottomrule
\end{tabular}
\end{center}

\vspace{0.6cm}
\begin{maginote}[Reporting a problem]
Include the run directory's \code{*\_config.json} and \code{*\_metadata.json},
the acceptance-harness JSON if you have it, and the seed. Most surprising
results trace back to a preprocessing contract mismatch that those files make
visible immediately.
\end{maginote}

% =====================================================================
\newpage
\bibliographystyle{unsrt}
% magi_shared.bib holds the entries also used by the paper; paper/ is
% gitignored, so they are copied here to keep docs/manual/ buildable from
% a fresh clone. magi_theory.bib holds the manual-only theory entries.
\bibliography{magi_shared,magi_theory}

\end{document}
